% 原书 PDF 171，印刷页 159。
\chapter{点的 Hilbert 概形}
\section*{引言}
在本章中，设 $k$ 是特征为 $0$ 的代数闭域。设 $X$ 为 $k$ 上的拟射影概形，带有丰沛线丛 $\mathcal O(1)$。$X$ 的 Hilbert 概形 $\operatorname{Hilb}(X)$ 参数化 $X$ 的所有闭且固有的子概形。我们知道，$\operatorname{Hilb}(X)$ 可以写成拟射影概形的不交并
\[
\operatorname{Hilb}(X)=\coprod_P\operatorname{Hilb}^P(X),
\]
其中 $\operatorname{Hilb}^P(X)$ 参数化 Hilbert 多项式为 $P$ 的子概形（即它表示这样的反变函子：将概形 $T$ 映到所有闭子概形 $Z\subset X\times T$ 所成的集合，其中 $Z$ 在 $T$ 上平坦，纤维的 Hilbert 多项式为 $P$）。见第二部分。

我们想处理最简单的情形，即 $P$ 为常数多项式 $n$。由于 Hilbert 多项式的次数就是子概形的维数，可见 $\operatorname{Hilb}^n(X)$ 参数化 $X$ 中长度为 $n$ 的零维子概形。换言之，这意味着
\[
\dim H^0(Z,\mathcal O_Z)
=\sum_{p\in\operatorname{supp}(Z)}\dim_k(\mathcal O_{Z,p})=n.
\]
这里 $\operatorname{len}(Z):=\dim H^0(Z,\mathcal O_Z)$ 为 $Z$ 作为自身上的模的长度。下文也将以 $X^{[n]}$ 表示 $\operatorname{Hilb}^n(X)$。$X^{[n]}$ 中元素的最简单例子，是由 $X$ 上 $n$ 个互异点组成的集合 $\{p_1,\ldots,p_n\}$。容易看出，它们构成 $X^{[n]}$ 的一个开子集。$X^{[n]}$ 参数化 $X$ 上 $n$ 个不必互异的点所成的集合，当其中一些点汇合时，还附带额外的非约化结构。另一个以不同方式参数化 $X$ 上 $n$ 点集的空间，是 $X$ 的对称幂 $X^{(n)}$，即 $X^n$ 对 $n$ 元对称群 $S_n$ 置换因子的作用的商。$X^{(n)}$ 参数化 $X$ 上次数为 $n$ 的有效零维圈，即形式和 $\sum_i n_i[p_i]$，其中 $p_i\in X$、$n_i\in\mathbb Z_{\geq0}$，且 $\sum n_i=n$。有一个显然的集合映射
\[
\rho:X^{[n]}\longrightarrow X^{(n)},\qquad
Z\longmapsto\sum_p\operatorname{len}(\mathcal O_{Z,p})[p],
\]
它将 $Z$ 映到带重数的支集。下面将会看到，这确实是概形的态射。

当 $X$ 为非奇异曲线时，我们将证明，对所有 $n$，$\rho$ 都是同构。若 $X$ 为非奇异曲面，则 $X^{[n]}$ 非奇异且不可约，$\rho$ 是 $X^{(n)}$ 的双有理奇点消解。这对于维数至少为 $3$ 的 $X$ 并不成立。
% 原书 PDF 172，印刷页 160。

曲面上点的 Hilbert 概形在数学和理论物理中都受到广泛关注。这部分是因为它们给出对称幂的典范奇点消解，也因为它们与向量丛的模空间及无限维 Lie 代数有关。关于其中若干联系的综述，见 [Göt02]。关于点的 Hilbert 概形及一些较新结果，既优美又易读的介绍是 [Nak99]。

\section{对称幂与 Hilbert--Chow 态射}
如引言所述，$X$ 上 $n$ 点的 Hilbert 概形 $X^{[n]}$ 与对称幂 $X^{(n)}=X^n/S_n$ 密切相关。我们首先需要知道，对称幂作为代数簇存在。

\begin{origstatement}{定义}{7.1.1}
设 $X$ 为 $k$ 上的拟射影簇，群 $G$（通过自同构）作用于 $X$。若下列条件成立，则称簇 $Y$ 连同满态射 $\pi:X\to Y$ 为 $X$ 对 $G$ 的\emph{商}（quotient）。
\begin{enumerate}
\item $\pi$ 的纤维是 $G$ 的轨道。
\item 到概形 $Z$ 的任一 $G$ 不变态射 $\varphi:X\to Z$ 都通过 $\pi$ 分解。
\end{enumerate}
由此可知，商若存在，则唯一到同构。以 $X/G$ 表示它。
\end{origstatement}
一般地，商是否存在是一个困难问题。不过，当 $G$ 有限且 $X$ 拟射影时，问题很容易。

\begin{origstatement}{定理}{7.1.2}
设 $X$ 为带有有限群 $G$ 作用的拟射影簇。则商 $X/G$ 作为簇存在。
\end{origstatement}
\begin{proof}[证明（梗概）]
先假定 $X$ 仿射。以 $k[X]$ 表示仿射坐标环。商的定义中的条件 (2) 意味着，$k[X/G]$ 应当是不变元环 $k[X]^G\subset k[X]$。容易看出，$k[X]^G$ 是有限生成 $k$ 代数，因此定义 $X/G:=\Spec(k[X]^G)$，这是一个仿射簇。令 $\pi:X\to X/G$ 为包含映射 $k[X]^G\subset k[X]$ 所诱导的态射。不难证明，$\pi$ 为满射，纤维为 $G$ 轨道。

若 $X$ 不仿射，则它具有仿射开覆盖 $(U_i)$；由于 $G$ 有限，可以选取这些仿射开集，使每条轨道都包含于某个 $U_i$。以 $W_i=\bigcap_{g\in G}g(U_i)$ 替代 $U_i$，便得到由 $G$ 不变仿射开子集组成的 $X$ 的开覆盖（因为 $X$ 是拟射影簇，从而分离，所以仿射开集的交仍仿射）。这时不难证明，$W_i/G$ 可粘合成商 $X/G$。
\end{proof}
进一步仔细论证（例如见 [Har92] 第 10 讲），可证明：若 $G$ 有限且 $X$ 拟射影，则 $X/G$ 拟射影；若 $X$ 射影，则 $X/G$ 也射影。

特别地，若 $X$ 为拟射影簇，则对称幂 $X^{(n)}:=X^n/S_n$（其中对称群 $S_n$ 通过置换因子作用）作为拟射影簇存在。

\begin{origstatement}{例}{7.1.3}
\begin{enumerate}
\item 由对称函数基本定理，$k[x_1,\ldots,x_n]^{S_n}=k[s_1,\ldots,s_n]$，其中 $s_i$ 为变量 $x_i$ 的初等对称函数。因此 $(\mathbb A^1)^{(n)}=\mathbb A^n$。类似地，可证明 $(\mathbb P^1)^{(n)}=\mathbb P^n$。
% 原书 PDF 173，印刷页 161。
\item 我们要说明，$(\mathbb A^2)^{(2)}$ 同构于 $\mathbb A^2$ 与二次锥面 $\Spec(k[u,v,w]/(uw-v^2))$ 的乘积。设 $x_1,y_1$ 和 $x_2,y_2$ 为两个因子的坐标。令 $x:=x_1-x_2$、$y=y_1-y_2$。则 $x,y,x':=x_1+x_2,y':=y_1+y_2$ 也是 $(\mathbb A^2)^2$ 上的坐标，$S_2$ 中的对换 $\tau$ 按 $\tau(x)=-x$、$\tau(y)=-y$、$\tau(x')=x'$、$\tau(y')=y'$ 作用。因此 $k[x,y,x',y']^\tau=(k[x,y]^\tau)[x',y']$。注意，$x,y$ 的每个单项式都是 $\tau$ 的特征向量，因为 $\tau(x^iy^j)=(-1)^{i+j}x^iy^j$；即 $k[x,y]$ 有一组由 $\tau$ 的特征向量组成的基。因此，$k[x,y]^\tau$ 是由特征值为 $1$ 的特征向量 $x^iy^j$ 张成的线性子空间，即要求 $i+j$ 为偶数。令 $u=x^2$、$v=xy$、$w=y^2$，可见 $k[x,y,x',y']^{S_2}$ 是由 $u,v,w,x',y'$ 生成的子代数，它同构于 $k[u,v,w,x',y']/(uw-v^2)$。特别地，可见 $(\mathbb A^2)^{(2)}$ 的奇异轨迹是 $(\mathbb A^2)^2$ 中对角线的像。

更一般地，若 $S$ 为非奇异曲面，则 $S$ 在平展拓扑下局部同构于 $\mathbb A^2$；因此，同样的论证表明，$S^{(2)}$ 的奇异轨迹是对角线的像。若 $k=\mathbb C$，则在解析拓扑中也同样成立。
\end{enumerate}
\end{origstatement}
按定义，对称幂 $X^{(n)}$ 的点是 $X$ 上的 $n$ 元点组在置换作用下的轨道，即有效零维圈 $\sum n_i[x_i]$，其中 $x_i\in X$、$n_i>0$，且 $\sum n_i=n$。由此可将 $X^{(n)}$ 另行描述为零维圈的 Chow 簇。

设 $X\subset\mathbb P^d$ 为局部闭的。可见 $X^{(n)}$ 是 $(\mathbb P^d)^{(n)}$ 的局部闭子簇。以 $\check{\mathbb P}^d$ 表示 $\mathbb P^d$ 的超平面所成的对偶射影空间。以 $\operatorname{Div}^n(\check{\mathbb P}^d)\simeq\mathbb P^{\binom{n+d}{d}-1}$ 表示 $\check{\mathbb P}^d$ 上次数为 $n$ 的有效除子的空间。对任意 $p\in\mathbb P^d$，令
\[
H_p:=\{l\in\check{\mathbb P}^d\mid p\in l\}.
\]
于是 $p\mapsto H_p$ 定义了同构 $\mathbb P^d\simeq\operatorname{Div}^1(\check{\mathbb P}^d)$。对 $(x_1,\ldots,x_n)\in(\mathbb P^d)^n$，令
\[
ch(x_1,\ldots,x_n):=\sum_i H_{x_i}\in\operatorname{Div}^n(\check{\mathbb P}^d).
\]
则 $ch:(\mathbb P^d)^n\to\operatorname{Div}^n(\check{\mathbb P}^d)$ 是 $S_n$ 不变态射，因而给出态射 $ch:(\mathbb P^d)^{(n)}\to\operatorname{Div}^n(\check{\mathbb P}^d)$。由于 $(\mathbb P^d)^n$ 射影，可见其像是闭的。可以验证，$ch$ 到其像为同构。特别地，也可以将 $X^{(n)}$ 与它在 $\operatorname{Div}^n(\check{\mathbb P}^d)$ 中的像等同。

现在定义 Hilbert--Chow 态射 $\rho:X^{[n]}_{\mathrm{red}}\to X^{(n)}$。我们将定义一个像的支集包含于 $X^{(n)}$ 的态射 $X^{[n]}\to\operatorname{Div}^n(\check{\mathbb P}^d)$ 来给出它。$\operatorname{Div}^n(\check{\mathbb P}^d)$ 表示这样的反变函子：为每个概形 $T$ 配备 $\check{\mathbb P}^d\times T$ 中次数为 $n$ 的有效相对 Cartier 除子 $D$。这里“相对的”有两个等价的含义：$D$ 在 $T$ 上平坦，或它在每个点 $t\in T$ 上的纤维上的限制都是 Cartier 除子。因此，为构造 $\rho$，需要一种得到有效 Cartier 除子的方法。

为此，先回顾 Mumford 的一个构造：在适当条件下，它为概形 $Y$ 上的凝聚层 $\mathcal F$ 配备 $Y$ 上的有效 Cartier 除子 $\operatorname{div}(\mathcal F)$。我们不会给出这一构造的全部细节或最一般形式，而只在关心的情形下给出梗概。一般构造可见 [MFK94] 第 5 章第 3 节。
% 原书 PDF 174，印刷页 162。

\begin{origstatement}{引理}{7.1.4}
设 $R$ 为局部环，$M$ 为 $R$ 模，具有解消
\[
0\longrightarrow R^n\xrightarrow{\varphi}R^n\longrightarrow M\longrightarrow0.
\]
则 $\det\varphi\in R/R^*$ 的类只依赖于 $M$，不依赖于所选的解消（即 $\det\varphi$ 在相差 $R$ 中一个单位的乘法的意义下定义良好）；此外，$\det\varphi$ 不是零因子。
\end{origstatement}
\begin{proof}
固定映射 $R^n\xrightarrow{\psi}M\to0$，所有可能的 $\varphi':R^n\to R^n$ 都由 $\varphi$ 与 $R^n$ 的一个自同构 $\alpha$ 复合得到。因此 $\det(\varphi')=\det(\alpha)\det(\varphi)$，且 $\det(\alpha)\in R^*$。

映射 $\psi:R^n\to M$ 对应于选取 $M$ 作为 $R$ 模的一组生成元。逐次添加和删除生成元，可以得到其他任何一种选择。因此，设 $\psi:R^n\to M$ 由 $m_1,\ldots,m_n\in M$ 给出，$\psi':R^{n+1}\to M$ 由 $m_1,\ldots,m_n,x\in M$ 给出。则在 $M$ 中有关系
\[
\sum_{i=1}^n a_im_i+x=0,\qquad a_i\in R;
\]
给定解消 $0\to R^n\xrightarrow{\varphi}R^n\xrightarrow{\psi}M\to0$，便得到解消
\[
0\longrightarrow R^{n+1}\xrightarrow{\varphi'}R^{n+1}
\xrightarrow{\psi'}M\longrightarrow0,\qquad
\varphi'=\begin{pmatrix}\varphi&0\\a_1\ \cdots\ a_n&1\end{pmatrix},
\]
且 $\det(\varphi')=\det(\varphi)$。
\end{proof}
\begin{origstatement}{注}{7.1.5}
设 $X$ 为 $k$ 上的光滑连通簇，$\mathcal F$ 为 $X$ 上满足 $\operatorname{supp}(\mathcal F)\ne X$ 的凝聚层。对不可约超曲面 $V\subset X$，以 $[V]$ 表示其泛点，并令 $R:=\mathcal O_{X,[V]}$。设 $M$ 为 $\mathcal F$ 在 $[V]$ 处的茎。则 $R$ 为离散赋值环，由 [Har77] III.6.11A、III.6.12A，$M$ 的同调维数为 $1$，即存在自由解消
\[
0\longrightarrow R^n\xrightarrow{\varphi}R^n\longrightarrow M\longrightarrow0.
\]
注意，由于 $\operatorname{supp}(\mathcal F)\ne X$，左侧两个模的秩相同。令 $m_V\in\mathbb Z_{\geq0}$ 为 $\det(\varphi)\in R$ 的赋值。

从更几何的角度，可以如下描述：存在开子集 $U\subset X$，其与 $V$ 的交为开集，并在 $U$ 上有解消
\[
0\longrightarrow\mathcal O^n\xrightarrow{\varphi}\mathcal O^n\longrightarrow\mathcal F\longrightarrow0.
\]
于是，令 $m_V$ 为 $\det(\varphi)$ 在 $V$ 的一个开子集上的消失阶数。

由上一个引理，$m_V$ 不依赖于解消的选择。
\end{origstatement}
\begin{origstatement}{定义}{7.1.6}
在注 7.1.5 的假设下，令
\[
\operatorname{div}(\mathcal F):=\sum_V m_VV.
\]
按定义，这是 $X$ 上的有效 Cartier 除子。因为只有当 $V\subset\operatorname{supp}(\mathcal F)$ 时，$m_V$ 才可能非零，所以此和为有限和。
\end{origstatement}
我们想在相对情形下构造这样的除子 $\operatorname{div}(\mathcal F)$。在陈述和证明为此所需的主要定理之前，先汇集几个有用的结果。
% 原书 PDF 175，印刷页 163。

\begin{origstatement}{引理}{7.1.7}
设 $X$ 为 $n$ 维光滑射影簇，$\mathcal F$ 为 $X\times S$ 上在 $S$ 上平坦的凝聚层。则 $\mathcal F$ 具有长为 $n$ 的局部自由解消。
\end{origstatement}
\begin{proof}
第一步：证明 $\mathcal F$ 是一个局部自由层的商。

对每个 $s\in S$，以 $\mathcal F_s$ 表示 $\mathcal F|_{X\times\{s\}}$。由于 $X$ 射影且 $S$ 为有限型，存在 $N>0$，使 $\mathcal F_s(N)$ 由整体截面生成，且对所有 $i>0$ 和所有 $s\in S$，都有 $H^i(\mathcal F_s(N))=0$。由平坦性，$\pi_*\mathcal F(N)$ 局部自由，其中 $\pi:X\times S\to S$ 为投影。因此得到满射 $\pi^*\pi_*\mathcal F(N)\to\mathcal F(N)$；与 $\mathcal O(-N)$ 作张量积，得到满射 $\mathcal E_0\to\mathcal F$，其中 $\mathcal E_0:=(\pi^*\pi_*\mathcal F(N))\otimes\mathcal O(-N)$ 局部自由。

第二步：存在 $X\times S$ 上凝聚层的正合列
\[
0\longrightarrow\mathcal E_n\longrightarrow\mathcal E_{n-1}
\longrightarrow\cdots\longrightarrow\mathcal E_0\longrightarrow\mathcal F\longrightarrow0,
\]
其中 $i=0,\ldots,n-1$ 时每个 $\mathcal E_i$ 都局部自由，且 $\mathcal E_n$ 在 $S$ 上平坦。

令 $\mathcal E_0\to\mathcal F$ 为第一步构造的满射；注意，它的核在 $S$ 上平坦。对这个核重复第一步，并继续归纳即可。

最后的第三步：证明 $\mathcal E_n$ 局部自由。

对每个 $s\in S$，$\mathcal E_n$ 在 $X\times\{s\}$ 上的限制都局部自由（由 [Har77] 命题 III.6.12A）。这一断言是局部的：因此，为完成证明，只需使用下面的引理。
\end{proof}
\begin{origstatement}{引理}{7.1.8}
设 $R$ 为正则局部环，$A$ 为 $(Loc/k)$ 中任意局部环（见定义 6.1.12）。设 $M$ 为有限生成 $S:=A\otimes_k R$ 模，在 $A$ 上平坦，并假定 $M\otimes_S R$ 作为 $R$ 模自由。则 $M$ 作为 $S$ 模自由。
\end{origstatement}
\begin{proof}
注意 $M\otimes_S R=M\otimes_A k$。设 $\bar m_1,\ldots,\bar m_r$ 为 $M\otimes_S R$ 作为 $R$ 模的一组自由生成元。将它们提升为 $M$ 中的元素 $m_1,\ldots,m_r$，并考虑由 $\alpha(e_i)=m_i$ 定义的同态 $\alpha:S^{\oplus r}\to M$；由 Nakayama 引理，它为满射。考虑所诱导的正合列
\[
0\longrightarrow K\longrightarrow S^{\oplus r}\xrightarrow{\alpha}M\longrightarrow0;
\]
由于 $M$ 在 $A$ 上平坦，该序列在 $\otimes_A k$ 后仍正合，这意味着 $K\otimes_A k=K\otimes_S R$ 为零。再次由 Nakayama 引理，可知 $K$ 为零。
\end{proof}
在本节余下部分，我们也考虑非闭点。

\begin{origstatement}{定义}{7.1.9}
设 $Y$ 为概形。若 $\operatorname{depth}\mathcal O_{Y,p}=d$，则称 $Y$ 的点 $p$（不必为闭点）的\emph{深度}（depth）为 $d$。特别地，它的深度为 $1$，是指 $m_p\subset\mathcal O_{Y,p}$ 含有一个非零因子，且对每个非零因子 $t\in m_p$，商环 $\mathcal O_{Y,p}/(t)$ 的极大理想中的每个元素都是零因子。
\end{origstatement}
\begin{origstatement}{注}{7.1.10}
若 $Y$ 为光滑簇，则深度为 $1$ 的点就是素除子即闭不可约超曲面的泛点。
\end{origstatement}
下述引理的证明可见 [EGAIV4] 命题 21.1.8。

\begin{origstatement}{引理}{7.1.11}
设 $Y$ 为概形，$f,g$ 为 $Y$ 上两个 Cartier 除子（即 $\mathcal K^*/\mathcal O^*$ 的整体截面）。则 $f=g$，当且仅当它们在每个深度为 $1$ 的点处的茎相同。
\end{origstatement}
% 原书 PDF 176，印刷页 164。
\begin{origstatement}{引理}{7.1.12}
(1) 设 $\mathcal E$ 为概形 $Y$ 上秩为 $r$ 的自由凝聚层。则 $\det\mathcal E:=\bigwedge^r\mathcal E$ 同构于 $\mathcal O_Y$；此同构在相差 $\mathcal O_Y^*$ 的一个截面的意义下唯一。

(2) 设
\[
0\longrightarrow\mathcal E_n\longrightarrow\mathcal E_{n-1}
\longrightarrow\cdots\longrightarrow\mathcal E_0\longrightarrow0
\]
为概形 $Y$ 上局部自由凝聚层的正合列。则有典范同构
\[
\bigotimes_{i=0}^n(\det\mathcal E_i)^{(-1)^i}\longrightarrow\mathcal O_Y.
\]
\end{origstatement}
\begin{proof}
(1) 是平凡的。(2) 用归纳法证明。先考虑 $n=2$ 的情形。若 $0\to\mathcal E_2\to\mathcal E_1\to\mathcal E_0\to0$ 为局部自由层的短正合列，则有典范同构
\[
\det\mathcal E_2\otimes\det\mathcal E_0\simeq\det\mathcal E_1.
\]
事实上，在一个开集上局部选取 $\mathcal E_1$ 的一组基 $f_1,\ldots,f_{r_2},g_1,\ldots,g_{r_0}$，使 $f_i$ 为 $\mathcal E_2$ 的一组基的像，而 $g_i$ 映到 $\mathcal E_0$ 的一组基（均以相同字母表示）。则同构由下式给出：
\[
(f_1\wedge\cdots\wedge f_{r_2})\otimes(g_1\wedge\cdots\wedge g_{r_0})
\longmapsto f_1\wedge\cdots\wedge f_{r_2}\wedge g_1\wedge\cdots\wedge g_{r_0}.
\]
容易验证，它不依赖于 $g_i$ 的选择，因此粘合成一个整体同构。

对于归纳步骤，注意可以将给定的正合列分解成两个局部自由层的正合列
\[
\begin{gathered}
0\longrightarrow\widetilde{\mathcal E}\longrightarrow\mathcal E_1\longrightarrow\mathcal E_0\longrightarrow0,\\
0\longrightarrow\mathcal E_n\longrightarrow\cdots\longrightarrow\mathcal E_2\longrightarrow\widetilde{\mathcal E}\longrightarrow0,
\end{gathered}
\]
从而由归纳假设得到典范同构 $\bigotimes_{i=0}^n(\det\mathcal E_i)^{(-1)^i}\simeq\mathcal O_Y$。
\end{proof}
\begin{origstatement}{定理}{7.1.13}
设 $X$ 为光滑不可约射影簇，$S$ 为概形，$\mathcal F$ 为 $X\times S$ 上在 $S$ 上平坦的凝聚层。假定对所有 $s\in S$，有 $\operatorname{supp}(\mathcal F_s)\ne X$。则存在 $X\times S$ 上的有效 Cartier 除子 $\operatorname{div}(\mathcal F)$，满足：
\begin{enumerate}
\item $\operatorname{div}(\mathcal F)$ 的形成与基变换交换。
\item 若 $S$ 为一点，则 $\operatorname{div}(\mathcal F)$ 与定义 7.1.6 中的除子相同。
\end{enumerate}
\end{origstatement}
\begin{proof}
设
\begin{equation}
0\longrightarrow\mathcal E_n\longrightarrow\cdots\longrightarrow\mathcal E_0\longrightarrow\mathcal F\longrightarrow0\tag{7.1.1}
\end{equation}
为 $\mathcal F$ 在 $X\times S$ 上的有限局部自由解消；由引理 7.1.7，它存在。设 $U\subset X\times S$ 为使各 $\mathcal E_i$ 在其上自由的开子集。令 $V:=U\setminus\operatorname{supp}(\mathcal F)$。则由引理 7.1.12(2)，在 $V$ 上有典范同构
\[
\mathcal O_{X\times S}|_V\simeq
\left.\bigotimes_{i=0}^n(\det\mathcal E_i)^{(-1)^i}\right|_V.
\]
另一方面，由于各 $\mathcal E_i$ 自由，由引理 7.1.12(1)，在 $U$ 上有同构
\[
\bigotimes_{i=0}^n(\det\mathcal E_i)^{(-1)^i}\simeq\mathcal O_{X\times S},
\]
% 原书 PDF 177，印刷页 165；定理7.1.13证明续。
它在相差一个单位的意义下唯一。这些同构的复合定义非零截面 $f\in\mathcal O_X^*(V)$，它在相差 $\mathcal O_X^*(U)$ 的一个截面的意义下唯一。由于 $V$ 为 $U$ 的稠密开子集，$\mathcal O_X^*(V)$ 的任一截面都定义 $\mathcal K^*(U)$ 的一个截面；因此，选定解消 $\mathcal E_i$ 和开子集 $U$，便可配备 $U$ 上的一个 Cartier 除子。现在证明，如此局部构造的 Cartier 除子只依赖于 $\mathcal F$，不依赖于解消的选择；这将意味着，它们粘合成一个整体 Cartier 除子。

为验证 Cartier 除子定义良好，设 $v\in X\times S$ 为深度 $1$ 的点，并令 $V$ 为 $X\times S$ 中以 $v$ 为泛点的闭子集。设 $\widetilde{\mathcal E}$ 为 $\mathcal E_0\to\mathcal F\to0$ 的核，于是有 $S$ 平坦层的正合列
\[
0\longrightarrow\widetilde{\mathcal E}\xrightarrow{\varphi}\mathcal E_0\longrightarrow\mathcal F\longrightarrow0.
\]
取 $v$ 处的茎后，此序列仍然正合；它成为有限生成 $\mathcal O_{X\times S,v}$ 模的序列，其中 $\mathcal E_0$ 自由。由于局部环 $\mathcal O_{X\times S,v}$ 的深度为 $1$，模 $\widetilde{\mathcal E}_v$ 也必为自由模；因此可在 $X\times S$ 中选取 $v$ 的一个开邻域 $W$（即与 $V$ 的交非空的开子集），使得 $\mathcal E_i|_W$ 和 $\widetilde{\mathcal E}|_W$ 均自由。

在 $W$ 上，有典范同构
\[
\det\widetilde{\mathcal E}\simeq\bigotimes_{i=1}^n(\det\mathcal E_i)^{(-1)^{i-1}},
\]
且 $\operatorname{div}(\mathcal F)$ 在 $W$ 上与 $\det\varphi$ 一致（如引理 7.1.4 中的构造）。

由引理 7.1.11，两个 Cartier 除子若在深度为 $1$ 的点处一致就相等；因此 $(f)$ 不依赖于解消的选择，从而粘合成 $X\times S$ 上的有效 Cartier 除子 $\operatorname{div}(\mathcal F)$。

此外，由于光滑簇上深度为 $1$ 的点恰好是素除子的泛点，这也说明，当 $S$ 为一点时，所得定义与定义 7.1.6 相同。

设 $h:T\to S$ 为态射。对 $X\times S$ 上的层 $\mathcal E$，以 $\mathcal E_T$ 表示其经 $id_X\times h$ 的拉回。利用 $\mathcal F$ 的平坦性，可以验证，若 $0\to\mathcal E_n\to\cdots\to\mathcal E_0\to\mathcal F\to0$ 为 $\mathcal F$ 在 $X\times S$ 上的解消，则
\[
0\longrightarrow(\mathcal E_n)_T\longrightarrow\cdots
\longrightarrow(\mathcal E_0)_T\longrightarrow(\mathcal F)_T\longrightarrow0
\]
为 $(\mathcal F)_T$ 的解消。因此，$\operatorname{div}(\mathcal F)$ 的拉回是 $\operatorname{div}(\mathcal F_T)$，即 $\operatorname{div}(\mathcal F)$ 与基变换相容。
\end{proof}
最后，可以构造 Hilbert--Chow 态射了。设 $H\subset\mathbb P^d\times\check{\mathbb P}^d$ 为关联对应。$H$ 是 $\mathbb P^d$ 上以 $\mathbb P^{d-1}$ 为纤维的纤维丛。以 $p,\check p$ 分别表示到 $\mathbb P^d$ 和 $\check{\mathbb P}^d$ 的投影。设 $S$ 为概形，$Z\subset\mathbb P^d\times S$ 为在 $S$ 上次数为 $n$ 的平坦闭子概形。令 $p_S:=p\times id_S$、$\check p_S:=\check p\times id_S$，$Z^*:=p_S^{-1}(Z)\subset H\times S$。令 $\mathcal F:=(\check p_S)_*(\mathcal O_{Z^*})$。则 $\mathcal F$ 是 $\check{\mathbb P}^d\times S$ 上在 $S$ 上平坦的凝聚层。设 $Z_s$ 为 $Z$ 在 $s\in S$ 上的纤维。则
\[
\operatorname{supp}(\mathcal F_s)=\{l\in\check{\mathbb P}^d\mid l\cap Z_s\ne\varnothing\},
\]
特别地，$\operatorname{supp}(\mathcal F_s)\ne\check{\mathbb P}^d$。因此，$\operatorname{div}(\mathcal F)$ 是 $\check{\mathbb P}^d\times S$ 上的相对 Cartier 除子。这样便构造了态射 $\rho:X^{[n]}\to\operatorname{Div}^n(\check{\mathbb P}^d)$。最后，直接由定义可以验证，$\operatorname{div}(\mathcal F_s)=\sum_{p\in\operatorname{supp}(Z_s)}\operatorname{len}(\mathcal O_{Z,p})H_p$。因而可见，
% 原书 PDF 178，印刷页 166；上段续。
$X^{[n]}$ 的像的支集为 $X^{(n)}$，所以若赋予 $X^{[n]}$ 以约化结构，该态射便通过 $X^{(n)}$ 分解。于是我们证明了下列定理。

\begin{origstatement}{定理}{7.1.14}
设 $X$ 为光滑拟射影簇。存在满态射 $\rho:X^{[n]}_{\mathrm{red}}\to X^{(n)}$，它在点的层面上由 $Z\mapsto\sum_{p\in\operatorname{supp}(Z)}\operatorname{len}(\mathcal O_{Z,p})[p]$ 给出。
\end{origstatement}
\section{不可约性与非奇异性}
我们将证明，若 $X$ 为 $k$ 上的非奇异拟射影曲线或曲面，则 $X^{[n]}$ 不可约且非奇异。下面引理和定理的证明，部分受到 [Leh04] 的启发。

\begin{origstatement}{引理}{7.2.1}
设 $X$ 为 $k$ 上的连通簇。则对所有 $n\geq0$，$X^{[n]}$ 都连通。
\end{origstatement}
\begin{proof}
先回顾，Quot 概形使我们能够定义 $X$ 上任意凝聚层的射影化。

对 $X$ 上凝聚层 $\mathcal F$，令 $\mathbb P(\mathcal F):=\operatorname{Quot}^1(\mathcal F)$，则 $\mathbb P(\mathcal F)$ 参数化 $\mathcal F$ 的纤维的一维商。$\mathbb P(\mathcal F)$ 是带有到 $X$ 的态射的拟射影概形，且 $\mathbb P(\mathcal F)$ 在 $x\in X$ 上的纤维为
\[
\mathbb P(\{\lambda:\mathcal F(x)\to k(x)\text{ 为满射}\})\simeq\mathbb P(\mathcal F(x)).
\]
现在对 $n$ 归纳证明断言。$X^{[0]}$ 是对应于空集的一个点。假定已经证明 $X^{[n]}$ 连通。则 $X\times X^{[n]}$ 连通。设 $Z_n(X)\subset X\times X^{[n]}$ 为万有族，理想层为 $\mathcal I_{Z_n(X)}$。令 $\mathbb P:=\mathbb P(\mathcal I_{Z_n(X)})$，其投影为 $\pi:\mathbb P\to X\times X^{[n]}$。$\mathbb P$ 在 $(x,Z)$ 上的纤维为射影空间，因而连通。所以 $\mathbb P$ 连通。

在 $X\times\mathbb P$ 上有万有正合列
\[
0\longrightarrow\mathcal I\longrightarrow
\mathcal I_{\pi^{-1}(Z_n(X))}\longrightarrow Q\longrightarrow0.
\]
这里 $Q$ 为万有商线丛，核 $\mathcal I$ 是 $X\times\mathbb P$ 上的理想层，定义子概形 $\mathcal Z\subset X\times\mathbb P$。上面的正合列给出正合列
\[
0\longrightarrow Q\longrightarrow\mathcal O_{\mathcal Z}
\longrightarrow\mathcal O_{\pi^{-1}Z_n(X)}\longrightarrow0.
\]
由于 $\mathcal O_{\pi^{-1}Z_n(X)}$ 在 $\mathbb P$ 上平坦且次数为 $n$，而 $Q$ 平坦且次数为 $1$，可见 $\mathcal Z$ 在 $\mathbb P$ 上平坦且次数为 $n+1$。于是有态射 $\psi:\mathbb P\to X^{[n+1]}$，它在点上将 $(\lambda:\mathcal I_Z\to k(x))$ 映到理想为 $\ker(\lambda)$ 的 $X$ 的子概形。我们要说明 $\psi$ 为满射。设 $W\in X^{[n+1]}$，$p\in\operatorname{supp}(W)$。在乘以 $p$ 处极大理想 $m$ 的作用的核中，选取元素 $f\in\mathcal O_W$。设 $Z\subset W$ 为理想 $(f)$ 所定义的子概形，则 $Z\in X^{[n]}$。令 $f=g_0,g_1,\ldots,g_k$ 为 $\mathcal I_Z/m\mathcal I_Z$ 的一组基，并定义 $\lambda:\mathcal I_Z\to k(p)$ 为 $\sum a_ig_i\mapsto a_0$。则 $\mathcal I_W=\ker(\lambda)$。所以 $\psi$ 为满射，故 $X^{[n+1]}$ 连通。
\end{proof}
\begin{origstatement}{注}{7.2.2}
若 $X$ 为光滑不可约曲面 $S$，则可以证明，$\mathbb P=\mathbb P(\mathcal I_{Z_n(S)})$ 正是 $S\times S^{[n]}$ 沿万有族 $Z_n(S)$ 的爆破 [Ell]、[ES98]。那里还证明了，$\mathbb P(\mathcal I_{Z_n(S)})$ 是关联对应
\[
S^{[n,n+1]}:=\{(Z,W)\in S^{[n]}\times S^{[n+1]}\mid Z\subset W\}.
\]
% 原书 PDF 179，印刷页 167；注7.2.2续。
这里 $Z\subset W$ 是指 $Z$ 为 $W$ 的子概形。常用 $S^{[n,n+1]}$ 来归纳研究点的 Hilbert 概形的性质。
\end{origstatement}
设 $X$ 为 $d$ 维非奇异拟射影簇。令 $X_0^n\subset X^n$ 为由各 $p_i$ 互异的点组 $(p_1,\ldots,p_n)$ 构成的稠密开集。以 $X_0^{(n)}$ 表示它在 $X^{(n)}$ 中的像，后者参数化各 $p_i$ 互异的有效零维圈 $\sum_i[p_i]$。这同样是稠密开集。由于 $S_n$ 在 $(X^n)_0$ 上自由作用，可见 $X_0^{(n)}$ 非奇异，维数为 $nd$。以 $X_0^{[n]}$ 表示它在 $X^{[n]}$ 中的逆像。可以验证，$X_0^{[n]}$ 的每一点处的切空间维数都为 $nd$，且 $\rho|_{X_0^{[n]}}$ 为同构。因此，$X^{[n]}$ 包含非奇异开子集，它同构于 $X^{(n)}$ 的一个开子集。当 $X$ 为曲线或曲面时，可以用这一点证明 $X^{[n]}$ 非奇异且不可约。

\begin{origstatement}{定理}{7.2.3}
\begin{enumerate}
\item 设 $C$ 为不可约非奇异拟射影曲线，$n\geq0$。则 $C^{[n]}$ 非奇异且不可约，维数为 $n$。
\item （Fogarty [Fog68]）设 $S$ 为不可约非奇异拟射影曲面，$n\geq0$。则 $S^{[n]}$ 非奇异且不可约，维数为 $2n$。
\end{enumerate}
\end{origstatement}
\begin{proof}
令 $X=C$ 或 $X=S$，并令 $d=\dim(X)$。由于 $X^{[n]}$ 连通，且包含一个 $nd$ 维非奇异开子集，只需证明：对所有 $[Z]\in X^{[n]}$，切空间 $T_{[Z]}X^{[n]}$ 的维数均为 $nd$。首先，这会说明 $X^{[n]}$ 在闭包 $\overline{X_0^{[n]}}$ 的每一点处都非奇异。若 $X^{[n]}$ 可约，则由连通性，必有另一个不可约分支与 $\overline{X_0^{[n]}}$ 相交，而交点必为 $X^{[n]}$ 的奇异点。

已知 $T_{[Z]}X^{[n]}=\Hom_{\mathcal O_X}(\mathcal I_Z,\mathcal O_Z)$。将 $\Hom(\mathord\bullet,\mathcal O_Z)$ 应用于 $0\to\mathcal I_Z\to\mathcal O_X\to\mathcal O_Z\to0$，得到
\[
\begin{aligned}
\Hom_{\mathcal O_X}(\mathcal O_Z,\mathcal O_Z)&\longrightarrow
\Hom_{\mathcal O_X}(\mathcal O_X,\mathcal O_Z)\\
&\longrightarrow\Hom_{\mathcal O_X}(\mathcal I_Z,\mathcal O_Z)
\longrightarrow\operatorname{Ext}^1_{\mathcal O_X}(\mathcal O_Z,\mathcal O_Z).
\end{aligned}
\]
第一个映射是同构 $k^n\to k^n$。因此 $\Hom(\mathcal I_Z,\mathcal O_Z)\subset\operatorname{Ext}^1(\mathcal O_Z,\mathcal O_Z)$，只需证明 $ext^1(\mathcal O_Z,\mathcal O_Z)\leq nd$。

对于曲线 $C$，有 $\Hom(\mathcal O_Z,\mathcal O_Z)=H^0(\mathcal O_Z)=k^n$；由 Serre 对偶性，$\operatorname{Ext}^1(\mathcal O_Z,\mathcal O_Z)=H^0(\mathcal O_Z\otimes K_C)^\vee=k^n$。

对于曲面 $S$，有 $\Hom(\mathcal O_Z,\mathcal O_Z)=H^0(\mathcal O_Z)=k^n$；由 Serre 对偶性，$\operatorname{Ext}^2(\mathcal O_Z,\mathcal O_Z)=H^0(\mathcal O_Z\otimes K_S)^\vee=k^n$。因此只需证明
\[
\chi(\mathcal O_Z,\mathcal O_Z)=\sum_{i=0}^2 ext^i(\mathcal O_Z,\mathcal O_Z)=0.
\]
% 原书此处的求和未印交错符号(-1)^i；保留并在审校记录中指出。
设 $0\to\mathcal E_l\to\cdots\to\mathcal E_0\to\mathcal O_Z\to0$ 为 $\mathcal O_Z$ 在 $S$ 上的局部自由解消。则 $\sum_i(-1)^i\operatorname{rk}(\mathcal E_i)=0$，且
\[
\chi(\mathcal O_Z,\mathcal O_Z)
=\sum_{i=0}^l(-1)^i\chi(\mathcal E_i,\mathcal O_Z)
=\sum_{i=0}^l(-1)^i n\cdot\operatorname{rk}(\mathcal E_i)=0.
\]
\end{proof}
% 原书 PDF 180，印刷页 168。
\begin{origstatement}{注}{7.2.4}
注意，在这个证明中，并没有证明障碍空间 $\operatorname{Ext}^1_{\mathcal O_X}(\mathcal I_Z,\mathcal O_Z)$ 消失；事实上，它通常不为零。
\end{origstatement}
\begin{origstatement}{注}{7.2.5}
设 $X$ 为非奇异簇。则当 $n\leq3$ 时，$X^{[n]}$ 非奇异。
\end{origstatement}
\begin{proof}
令 $d=\dim(X)$。只需证明，对所有 $[Z]\in X^{[n]}$，有 $hom_{\mathcal O_X}(\mathcal I_Z,\mathcal O_Z)\leq dn$。显然，
\[
\Hom_{\mathcal O_X}(\mathcal I_Z,\mathcal O_Z)
=\bigoplus_{p\in\operatorname{supp}(Z)}
\Hom_{\mathcal O_{X,p}}(\mathcal I_{Z,p},\mathcal O_{Z,p}).
\]
因此可以归结到 $\operatorname{supp}(Z)$ 为一个点 $p$ 的情形。设 $m$ 为 $p$ 处的极大理想。容易证明，在 $p$ 处存在局部参数 $x_1,\ldots,x_d$，使 $\mathcal O_{Z,p}$ 具有下列形式：
\[
\begin{cases}
k[x_1,\ldots,x_d]/m,&n=1,\\
k[x_1,\ldots,x_d]/(m^2+(x_2,\ldots,x_d)),&n=2,\\
k[x_1,\ldots,x_d]/(m^3+(x_2,\ldots,x_d)),&n=3,\\
\text{或 }k[x_1,\ldots,x_d]/(m^2+(x_3,\ldots,x_d)).&
\end{cases}
\]
在所有情形下，都容易验证 $\underline{Hom}_{\mathcal O_{X,p}}(\mathcal I_{Z,p},\mathcal O_{Z,p})=d\operatorname{len}(\mathcal O_{Z,p})$。
% 原书末式以带下划线的Hom等于数值书写，忠实保留。
\end{proof}
\begin{origstatement}{注}{7.2.6}
设 $X$ 非奇异且维数为 $3$。则 $X^{[4]}$ 奇异。
\end{origstatement}
\begin{proof}
设 $[Z]\in X^{[4]}$ 为点 $\mathcal O_Z=\mathcal O_p/m^2$。则
\[
\Hom_{\mathcal O_X}(\mathcal I_Z,\mathcal O_Z)
=\Hom_k(m^2/m^3,m/m^2)=k^{18}.
\]
所以维数大于 $dn=12$。
\end{proof}
\begin{origstatement}{注}{7.2.7}
设 $X$ 为非奇异簇。注意，长度至多为 $3$ 的零维子概形 $Z\subset X$ 的嵌入维数为 $2$，即它局部是一个非奇异曲面的子概形。事实上，不难看出，若 $Z\subset X^{[n]}$ 的嵌入维数至多为 $2$，则 $[Z]$ 是 $X$ 的非奇异点。
% 原书前句为embedding dimension 2，后句写Z⊂X^[n]和点属于X；忠实保留原文疑点。
\end{origstatement}
\begin{origstatement}{例}{7.2.8}
Iarrobino 的一个例子 [Iar72] 表明，若 $X$ 非奇异，维数 $d\geq3$，而且 $n$ 充分大，则 $X^{[n]}$ 可约。

其思想很简单：Iarrobino 证明，当 $n$ 大时，$X^{[n]}$ 中集中于一点 $p\in X$ 的子概形所成的闭子集 $\rho^{-1}(n[p])$ 的维数大于或等于 $nd$。另一方面，我们知道，由 $n$ 个互异点构成的轨迹 $X_0^{[n]}$ 在 $X^{[n]}$ 中是开集，维数为 $nd$。因此，其闭包为维数 $nd$ 的不可约分支，不可能包含 $\rho^{-1}(n[p])$。

现在验证关于 $\rho^{-1}(n[p])$ 维数的结果。设 $m$ 为 $\mathcal O_{X,p}$ 的极大理想。考虑形如 $I=q^{-1}(V)\subset\mathcal O_{X,p}$ 的向量子空间，其中 $V$ 为 $m^k/m^{k+1}$ 中余维为 $l$ 的向量子空间，$q:=\mathcal O_{X,p}\to\mathcal O_{X,p}/m^{k+1}$ 为商映射。注意，对上述任意 $I$，$m$ 中任意元素在 $I/m^{k+1}$ 上的乘法作用都为零。因此，$I$ 为 $\mathcal O_{X,p}$ 的理想，定义子概形 $Z\in\rho^{-1}(n[p])$，其中 $n:=\binom{k+d-1}{d}+l$。这样，便得到 $\rho^{-1}(n[p])$ 的一个闭子集，它同构于 $s:=\binom{k+d-1}{d-1}$ 维向量空间 $m^k/m^{k+1}$ 的 $l$ 维商所成的 Grassmann 簇，因此维数为 $l(s-l)$。
% 原书 PDF 181，印刷页 169；例7.2.8续。

若 $d\geq3$，$k$ 充分大，且 $l$ 接近 $s/2$，则容易看出 $l(s-l)\geq nd$。例如，当 $d=3$、$k=7$ 时，得 $s:=36$，且 $\binom{k+d-1}{d}:=84$。因此，例如取 $l=12$，得到 $n=96$，而 $l(s-l)=288=3n$。所以，对光滑三维簇 $X$，$X^{[96]}$ 可约。
\end{origstatement}
\begin{origstatement}{注}{7.2.9}
可以证明（见 [Iar85]）：
\begin{enumerate}
\item 对 $n\leq7$ 和所有 $d$，$(\mathbb P^d)^{[n]}$ 不可约；
\item 当 $d=3$ 时，$(\mathbb P^d)^{[8]}$ 不可约，而当 $d\geq4$ 时它可约。
\end{enumerate}
不难证明，将 $\mathbb P^d$ 换成任意 $d$ 维光滑拟射影簇 $X$，这些结果仍然成立。
\end{origstatement}
\section{Hilbert 概形的例子}
设 $X$ 为 $d$ 维非奇异射影簇。当 $n$ 较小时，我们给出几个 $X^{[n]}$ 的例子。

\begin{origstatement}{例}{7.3.1}
\begin{enumerate}
\item $X^{[0]}$ 是一个约化点，对应于 $X$ 的空子概形。
\item $X$ 的长度为 $1$ 的子概形就是 $X$ 的点，且 $X^{[1]}=X$。万有族就是对角线 $\Delta\subset X\times X$。
\item $X^{[2]}$ 的点要么是 $X$ 上两个互异点的集合 $\{p_1,p_2\}$，要么是集中于一点 $p$ 的长度为 $2$ 的子概形 $Z$。设 $m$ 为 $p$ 的理想，则 $m\supset\mathcal I_Z\supset m^2$。因此，$\mathcal I_Z$ 由 $m/m^2$ 的一个余维为 $1$ 的子空间给出，即由 $\mathbb P(T_{X,p})=\mathbb P^{n-1}$ 中的一点给出。换言之，$X^{[2]}$ 的一点要么是 $X$ 的两个点组成的集合，要么是一个点 $p$ 及 $p$ 处的一条切方向。
% 原书这里写P^{n-1}，本段维数记为d；保留并记录。

由此可如下整体地描述 $X^{[2]}$：设 $\widetilde{X^2}$ 为 $X\times X$ 沿对角线的爆破，$E$ 为例外除子。$S_2$ 在 $X^2$ 上的作用扩张到 $\widetilde{X^2}$。设 $Y$ 为商。$E$ 是 $S_2$ 的非平凡元素的不动点轨迹。由于它是除子，可见 $Y$ 非奇异。容易看出，$\pi:\widetilde{X^2}\to Y$ 平坦且次数为 $2$。这给出双有理且双射的态射 $Y\to X^{[2]}$。因此，由 Zariski 主定理，它为同构。
\end{enumerate}
\end{origstatement}
\begin{origstatement}{习题}{7.3.2}
验证 $\operatorname{Hilb}^2\mathbb A^2$ 是 $\mathbb A^2$ 的第二对称幂沿对角线的爆破。提示：显式构造一个万有族。事实上，将 $\mathbb A^2$ 换成任意光滑代数簇，这一点仍然成立。
\end{origstatement}
对一般的 $n$，若 $X$ 的维数为 $1$ 或 $2$，则仍可以描述 $X^{[n]}$ 的一些性质。对于非奇异曲线，点的 Hilbert 概形就是对称幂。

\begin{origstatement}{命题}{7.3.3}
设 $C$ 为非奇异拟射影曲线。则 $\rho:C^{[n]}\to C^{(n)}$ 为同构。
\end{origstatement}
\begin{proof}
由于 $C$ 在点 $p$ 处的局部环为离散赋值环，$\mathcal O_{C,p}$ 中所有理想都是极大理想 $m_p$ 的幂。因此，对所有 $[Z]\in C^{[n]}$，有
\[
\mathcal O_Z=\bigoplus_i\mathcal O_{C,p_i}/m_{p_i}^{n_i},\qquad
\sum_i n_i=n.
\]
% 原书 PDF 182，印刷页 170；命题7.3.3证明续。
于是，$\rho$ 将 $Z$ 映到 $\sum_i n_i[p_i]$。因此 $\rho$ 为双射。又因为它也是双有理的，由 Zariski 主定理，$\rho$ 为同构。

另一种方法是注意到
\[
\pi:C\times C^{(n-1)}\longrightarrow C^{(n)},\qquad
\left(p,\sum n_i[p_i]\right)\longmapsto\sum n_i[p_i]+[p]
\]
在 $C^{(n)}$ 上平坦且次数为 $n$，从而定义 $\rho$ 的逆。
\end{proof}
现在设 $S$ 为非奇异曲面。我们已经看到，$S^{(n)}$ 中由 $n$ 个互异点之和构成的开子集 $S_0^{(n)}$ 非奇异。现在说明，$S^{(n)}\setminus S_0^{(n)}$ 是 $S^{(n)}$ 的奇异轨迹。由例 7.1.3，当 $n=2$ 时，这一点成立，因为 $S^{(2)}\setminus S_0^{(2)}$ 是 $S^2$ 中的对角线在商映射 $S^2\to S^{(2)}$ 下的像。设 $\Delta=\bigcup_{1\leq i<j\leq n}\Delta_{i,j}$ 为 $S^n$ 中的大对角线，其中 $\Delta_{i,j}=\{(x_1,\ldots,x_n)\in S^n\mid x_i=x_j\}$，并令 $\Delta_0$ 为其中恰有两个 $x_i$ 相等的开子集。设 $p\in\Delta_0$；可以假定 $p=(z,z,z_3,\ldots,z_n)\in\Delta_{1,2}$。则 $p$ 在 $S_n$ 中的稳定子为 $\{1,\tau\}$，其中 $\tau$ 为交换第一、第二因子的对换。因此由例 7.1.3，$p$ 在 $S^{(n)}$ 中的像的一个形式邻域（即 $S^{(n)}$ 在 $p$ 处的完备局部环）同构于 $k[[u,v,w,x',y',x_3,y_3,\ldots,x_n,y_n]]/(uw-v^2)$。所以，$\Delta_0$ 的像中的所有点都是 $S^{(n)}$ 的奇异点。$\Delta_0$ 的像的闭包为 $S^{(n)}\setminus S_0^{(n)}$。由于 $S^{(n)}$ 的奇异轨迹是闭集，可知 $S^{(n)}$ 的奇异轨迹恰为 $S^{(n)}\setminus S_0^{(n)}$。

\begin{origstatement}{定理}{7.3.4}
（Fogarty [Fog68]）设 $S$ 为非奇异拟射影曲面。则 $\rho:S^{[n]}\to S^{(n)}$ 为奇点消解。
\end{origstatement}
\begin{proof}
$S^{[n]}$ 非奇异且不可约，而 $\rho$ 在开子集 $S_0^{(n)}$ 上为同构，该开子集正是 $S^{(n)}$ 的非奇异轨迹。因此，它为奇点消解。
\end{proof}
\begin{origstatement}{注}{7.3.5}
\begin{enumerate}
\item 曲面上点的 Hilbert 概形受到关注的原因之一，是它们给出对称幂的典范奇点消解。
\item 事实上，已经知道 [Bea83]，$\rho:S^{[n]}\to S^{(n)}$ 有一个特别好的性质：$S^{(n)}$ 为 Gorenstein，且 $K_{S^{[n]}}=\rho^*(K_{S^{(n)}})$。这时称 $\rho$ 为保典（crepant）奇点消解。
\item 当 $S$ 为 K3 曲面或 Abel 簇时，[Bea83] 还证明，$S^{[n]}$ 为全纯辛簇，即它具有处处非退化的全纯二形式。这是点的 Hilbert 概形受到关注的另一个原因，因为全纯辛簇非常稀少。事实上，几乎所有已知的全纯辛射影簇都可以与点的 Hilbert 概形联系起来。
% 原文homomorphic symplectic系明显文字误拼；依据同句全纯二形式定义译为全纯辛。
\end{enumerate}
\end{origstatement}
\section{Hilbert 概形的一个分层}
在第七章的余下部分，设 $S$ 为复数域上的光滑射影曲面。

我们要研究 $S^{(n)}$ 和 $S^{[n]}$ 的一个自然分层。
% 原书 PDF 183，印刷页 171。
\begin{origstatement}{定义}{7.4.1}
对 $n$ 的任意分拆 $\nu=(n_1,\ldots,n_r)$（即 $n_1\geq n_2\geq\cdots\geq n_r>0$，且 $\sum n_i=n$），定义 $S^{(n)}$ 的局部闭子集
\[
S_\nu^{(n)}:=\left\{\sum n_i[x_i]\in S^{(n)}\mid
x_i\in S\text{ 为互异点}\right\}.
\]
这样便有由局部闭层组成的分层
\[
S^{(n)}=\coprod_\nu S_\nu^{(n)}.
\]
令 $S_\nu^{[n]}:=\rho^{-1}(S_\nu^{(n)})$，也得到 $S^{[n]}$ 的局部闭分层。
\end{origstatement}
现在研究各层 $S_\nu^{(n)}$。首先注意，它们非奇异。设 $\nu=(n_1,\ldots,n_r)$。写作 $\nu=(1^{\alpha_1},2^{\alpha_2},\ldots,s^{\alpha_s})$，其中 $\alpha_i$ 为 $i$ 在 $(n_1,\ldots,n_r)$ 中出现的次数。设 $(S^{\alpha_1}\times\cdots\times S^{\alpha_s})_*\subset S^{\alpha_1}\times\cdots\times S^{\alpha_s}$ 为所有分量 $p_i\in S$ 均互异的开子集。则 $S_\nu^{(n)}$ 是 $(S^{\alpha_1}\times\cdots\times S^{\alpha_s})_*$ 对 $S_{\alpha_1}\times\cdots\times S_{\alpha_s}$ 作用的商，其中每个 $S_{\alpha_i}$ 置换 $S^{\alpha_i}$ 的各因子。由于此作用自由，$S_\nu^{(n)}$ 非奇异，且其维数为 $\alpha_1+\cdots+\alpha_s=r$。
% 原书维数确印r而非2r；与下一节的2r不同，保留并记录。

\begin{origstatement}{习题}{7.4.2}
设 $\overline{S_\nu^{(n)}}$ 为 $S_\nu^{(n)}$ 在 $S^{(n)}$ 中的闭包。令
\[
\pi_\nu:S^{(\alpha_1)}\times\cdots\times S^{(\alpha_s)}\longrightarrow S^{(n)},\qquad
(\xi_1,\ldots,\xi_s)\longmapsto\sum_{i=1}^s i\xi_i.
\]
设 $(S^{(\alpha_1)}\times\cdots\times S^{(\alpha_s)})_*$ 为 $(S^{\alpha_1}\times\cdots\times S^{\alpha_s})_*$ 对 $S_{\alpha_1}\times\cdots\times S_{\alpha_s}$ 的商。

验证，上面的论证表明，$\pi_\nu$ 在 $(S^{(\alpha_1)}\times\cdots\times S^{(\alpha_s)})_*$ 上的限制是到 $S_\nu^{(n)}$ 的同构。因此，$\pi_\nu$ 的像是闭包 $\overline{S_\nu^{(n)}}$。证明，$\pi_\nu:S^{(\alpha_1)}\times\cdots\times S^{(\alpha_s)}\to\overline{S_\nu^{(n)}}$ 是正规化。
\end{origstatement}
\begin{origstatement}{定义}{7.4.3}
令 $H_n:=\operatorname{Hilb}^n(\Spec(k[[x,y]]/(x,y)^n))$，它参数化 $k[[x,y]]$ 中余长度为 $n$ 的理想（注意，$k[[x,y]]$ 中余长度为 $n$ 的每个理想都包含 $(x,y)^n$）。$H_n$ 称为\emph{点支撑 Hilbert 概形}（punctual Hilbert schemes）。
\end{origstatement}
\begin{origstatement}{命题}{7.4.4}
设 $\nu:=(n_1,\ldots,n_r)$ 为 $n$ 的一个分拆。则经 Hilbert--Chow 态射，$S_\nu^{[n]}$ 为 $S_\nu^{(n)}$ 上以 $H_{n_1}\times\cdots\times H_{n_r}$ 为纤维的局部平凡纤维丛（对解析拓扑和平展拓扑而言）。
\end{origstatement}
\begin{proof}
先考察最坏的一层 $\rho:S_{(n)}^{[n]}\to S_{(n)}^{(n)}\simeq S$。对 $S$ 中一点 $p$，纤维 $\rho^{-1}(n[p])$ 是 $S$ 中长度为 $n$、支集为 $p$ 的子概形 $Z$ 所成的集合。任意这样的概形的理想都包含于 $m_p^n$。于是，在 $p$ 的邻域 $U$ 中选取全纯局部坐标 $x,y$，就确定同构
\[
\rho^{-1}(U)\simeq U\times H_n.
\]
因此，$\rho:S_{(n)}^{[n]}\to S$ 为解析拓扑下以 $H_n$ 为纤维的局部平凡纤维丛。进一步仔细论证，可以在这个论证中将局部坐标换成 $p$ 的一个 Zariski 开邻域中的局部参数（即 $x,y$，使 $dx,dy$ 在 $U$ 的每一点都张成余切空间）。因此，这甚至是 Zariski 拓扑下的局部平凡纤维丛。
% 原书此处为ideal contained in m_p^n，和定义7.4.3方向相反；忠实保留。

现在设 $\nu=(n_1,\ldots,n_r)=(1^{\alpha_1},2^{\alpha_2},\ldots,s^{\alpha_s})$ 为 $n$ 的分拆。令 $\xi:=n_1[p_1]+\cdots+n_r[p_r]\in S_\nu^{(n)}$。则纤维 $\rho^{-1}(\xi)$ 同构于 $H_{n_1}\times\cdots\times H_{n_r}$。
% 原书 PDF 184，印刷页 172；命题7.4.4证明续。
事实上，可以在 $S$ 中为各 $p_i$ 选取（解析拓扑下的）互不相交的开邻域 $U_i$；它们给出 $\xi$ 在 $S_\nu^{(n)}$ 中的开邻域 $U$，使得
\[
\rho^{-1}(U)\simeq U\times H_{n_1}\times\cdots\times H_{n_r}.
\]
因此，$\rho:S_\nu^{[n]}\to S_\nu^{(n)}$ 是解析拓扑下以 $H_{n_1}\times\cdots\times H_{n_r}$ 为纤维的局部平凡纤维丛。

再次仔细论证，可以证明更多：此丛在平展拓扑下局部平凡。有
\[
(S^{\alpha_1+\cdots+\alpha_s})_*\times_{S_\nu^{(n)}}S_\nu^{[n]}
=\left(S_{(n_1)}^{[n_1]}\times\cdots\times S_{(n_r)}^{[n_r]}\right)_*,
\]
其中右侧的 $*$ 表示子概形的支集互异。$(S^{\alpha_1+\cdots+\alpha_s})_*\to S_\nu^{(n)}$ 为平展态射，而 $\left(S_{(n_1)}^{[n_1]}\times\cdots\times S_{(n_r)}^{[n_r]}\right)_*\to S^r$ 是 Zariski 拓扑下的局部平凡纤维丛；因此 $S_\nu^{[n]}\to S_\nu^{(n)}$ 平展局部平凡。
\end{proof}
现在考察 $\rho$ 的纤维。由上面的讨论，只需考察点支撑 Hilbert 概形 $H_n$。它们已被广泛研究（例如见 [Iar77]、[Bri77]）。综述还可见 [Iar87]。

显然，$H_1$ 为对应于极大理想 $(x,y)$ 的一个点，而 $H_2=\mathbb P_1$，对应于由 $(x,y)/(x,y)^2$ 中余维为 $1$ 的线性子空间所给出的理想。

当 $n\geq3$ 时，需要根据概形 $Z$ 的嵌入维数分两种情形。若概形 $Z\in H_n$ 的嵌入维数为 $1$，则称它为\emph{曲线型}（curvilinear）。这意味着它局部位于一条光滑曲线上。以 $H_n^c\subset H_n$ 表示曲线型子概形构成的开子概形。

\begin{origstatement}{注}{7.4.5}
$H_n^c$ 是 $\mathbb P^1$ 上的局部平凡 $\mathbb A^{n-2}$ 丛。特别地，$H_n^c$ 不可约，维数为 $n-1$。
\end{origstatement}
\begin{proof}
有态射 $H_n^c\to\mathbb P^1$，将 $Z$ 映到它在 $(0,0)$ 处的切空间。设 $U\subset H_n^c$ 为 $\mathbb A^1=\mathbb P^1\setminus(x=0)$ 的逆像。则 $Z\in U$ 位于一条曲线 $C$ 上，后者由幂级数 $y=\sum_{i\geq1}a_ix^i$ 的图像给出。因此可写成
\[
I_Z=(y-a_1x-\cdots-a_{n-1}x^{n-1},x^n).
\]
所以，$U$ 同构于以 $a_1,\ldots,a_{n-1}$ 为坐标的 $\mathbb A^{n-1}$，到 $\mathbb A^1$ 的映射由 $(a_1,\ldots,a_{n-1})\mapsto a_1$ 给出。$\mathbb P^1\setminus(y=0)$ 的逆像上的论证类似。
\end{proof}
当 $n=3$ 时，可见唯一不是曲线型的子概形是理想为 $(x,y)^2$ 的概形，且 $H_3^c$ 在 $H_3$ 中稠密。事实上，一般地也成立。

\begin{origstatement}{定理}{7.4.6}
[Bri77] $H_n$ 不可约，且 $H_n^c$ 在 $H_n$ 中为稠密开集。
\end{origstatement}
[Ell]、[ES98] 中关于这一不可约性的另一个证明，利用了如下事实：可经关联簇 $S^{[n,n+1]}=\{(Z,W)\in S^{[n]}\times S^{[n+1]}\mid Z\subset W\}=\mathbb P(\mathcal I_{Z_n(S)})$ 对 $S^{[n]}$ 作归纳描述（见注 7.2.2）。
